\chapter{可信机器学习基础}
\label{chap:trust-foundations}

上一部分讨论可学习性：在明确假设空间、损失、抽样关系和比较基准后，有限样本泛化保证说明，由样本得到的风险判断能够在多大误差和失败概率下推广到相应总体。这解决的是从有限样本到既定总体的统计推广问题。

从可学习性走向真实部署中的可使用性，还要把推广对象和使用条件说清。需要判断推广的是模型的预测误差、概率覆盖，还是模型、决策规则与使用流程共同造成的行动损失；也要回答受影响者能否质询决定并获得复核，以及数据、权限或流程变化后由谁发现失效、限制使用和执行补救。

一个邮件分类器在测试集上的准确率达到97\%，是否就可以让它自动删除被判为垃圾邮件的内容？如果它只负责把可疑邮件放入一个可恢复的文件夹，少量误判可能可以接受；如果它直接删除求职通知或合同附件，同样的误判率就可能造成更大的损失。即使有限样本保证支持它在某个总体上的误判率结论，也不能单独决定应当赋予隔离还是删除权限。

\term{可信机器学习}（trustworthy machine learning）研究模型及其使用过程满足合理使用要求的条件和证据。这里的可信性是相对于任务、受影响者和使用范围而言的。一个系统可以在概率预测上可靠，却泄露训练信息；也可以抵御某种攻击，却对少数群体持续产生较高的错误率。因此，可信性需要拆解为具体主张，再分别检验，而不能由一个笼统的综合得分代替。

本部分把讨论范围限定为技术可信性：考察由模型、数据、算法和技术系统证据支持的主张，并为组织审批、法律合规和公共治理提供可复核的输入。是否部署、由谁承担责任以及什么风险可以接受，仍需要有权主体结合领域规则和受影响者作出决定；技术测试通过不能替代这些判断，本部分也不声称完整覆盖人工智能治理。

本章从使用情境建立可信主张，再沿三个层次检查它：系统属性或规范要求说明必须满足什么，证据能力说明已有材料能够支持多强的结论，威胁与生命周期治理条件说明谁能使结论失效、由谁决定以及何时处置。无论考察可靠性、鲁棒性、公平性、行为约束、可解释性、隐私还是安全，核心问题都是：在明确边界与剩余不确定性的前提下，已有证据可以支持系统承担多大范围的决策？

\section{技术可信性的范围与三个层次}
\label{sec:trust-context}

风险分析需要先追踪模型输出怎样经由决策产生后果，再明确由谁承担这些后果、结论适用于什么范围。前者确定损失的含义，后者限定损失评价所针对的人和环境；缺少任一项，单个性能数值都无法支持使用判断。

可信性分析的起点是系统属性或规范要求，它回答系统在何种情境下必须满足什么。相应说明要写清受影响对象、评价单位和适用范围，再分别提出可靠性、鲁棒性、公平性、对齐、可解释性、隐私与安全要求。

要求明确后，需要考察证据能力：现有材料能够支持多强的结论。为此要说明显式假设、理论量词和抽样方式，并结合压力测试、机制分析与人因评价，记录反证和剩余不确定性。

证据结论还必须附带威胁与生命周期治理条件。分析要继续追问谁或什么能够使主张失效，谁负责决定，以及何时触发监测、限制使用、回退和恢复。

这三个层次相互依赖，却不能合并成一项指标。没有要求，测试结果不知道要支持哪种行动；没有证据，要求只是愿望；没有威胁和生命周期条件，一次通过的测试也会在数据、模型、权限或使用方式改变后被错误沿用。下面先用邮件系统明确情境与危害，再分别展开要求、证据和生命周期条件。

\subsection{从预测结果到实际后果}

模型输出通常只是行动的一个输入。设输入为$\bm{x}$，模型给出的结果为$f(\bm{x})$，使用者或系统根据规则$d$选择行动$a=d(f(\bm{x}),\bm{x})$。例如，分类器输出垃圾邮件概率，而决策规则可以选择放行、隔离或转交人工。即使$f$不变，改变$d$也会改变系统造成的后果。

为了比较这些后果，记真实状态为$y$，行动损失为$L(a,y)$。相对于部署分布$\mathcal P$，系统的期望损失为
\begin{equation}
 R_{\mathrm{sys}}(f,d)
 =\mathbb E_{(\bm{x},y)\sim\mathcal P}
 L\bigl(d(f(\bm{x}),\bm{x}),y\bigr).
 \label{eq:trust-system-risk}
\end{equation}
这个量同时依赖预测、决策规则和数据分布。用零一损失评价分类器，只计算是否预测错误；用行动损失评价系统，还要区分错误造成的代价。两种评价都可以有用，但回答的问题不同。

\begin{example}[准确率与行动损失]
\label{ex:trust-mail-cost}
考虑一个人为构造的1000封邮件测试集。模型甲误拦10封正常邮件、漏过20封垃圾邮件；模型乙误拦2封正常邮件、漏过40封垃圾邮件。甲的准确率为97\%，乙为95.8\%。若每次误拦损失为10，每次漏过损失为1，则甲的总损失为$10\times10+20=120$，乙为$2\times10+40=60$。在这一损失设定下，准确率较低的乙反而更适合使用。
\end{example}

示例中的代价由教学设定给出，并非邮件任务的普遍取值。实际使用还可能需要考虑延误、申诉和恢复成本。若隔离后的邮件可以及时找回，损失可能小于永久删除；若用户不知道发生了误拦，可恢复性也未必转化为实际补救。这要求评价范围覆盖完整使用流程。

当模型输出概率时，决策规则可以由损失比较推导。设$p$确实表示给定输入为垃圾邮件的概率，误拦正常邮件的成本为$c_{\mathrm{fp}}>0$，放过垃圾邮件的成本为$c_{\mathrm{fn}}>0$，正确处理的成本暂设为零。隔离与放行的条件期望损失分别为$(1-p)c_{\mathrm{fp}}$和$pc_{\mathrm{fn}}$，因此选择隔离的条件是
\begin{equation}
 p>\frac{c_{\mathrm{fp}}}{c_{\mathrm{fp}}+c_{\mathrm{fn}}}.
 \label{eq:trust-cost-threshold}
\end{equation}
相等时两种行动在当前损失模型下等价。阈值不是固定的0.5：误拦越昂贵，支持隔离所需的概率就越高。这里的$p$是给定输入后的真实条件概率；把模型估计值代入，只能得到依赖估计质量的决策。即使若干概率区间中的平均预测与实际频率相符，也还不能断言每个输入的概率都准确。第~\ref{chap:trust-reliability}章将区分这些可靠性要求，并说明如何检验它们。

\begin{example}[增加复核选项后的决策]
\label{ex:trust-review-cost}
在一个教学设定中，取$c_{\mathrm{fp}}=10$、$c_{\mathrm{fn}}=1$。若人工复核可以无误确认标签，每次固定成本为0.2，则三项候选损失为$10(1-p)$、$p$和0.2。$p<0.2$时放行更合算，$p>0.98$时隔离更合算，$0.2<p<0.98$时选择复核；两个边界处分别有两种行动成本相同。这个区域来自行动成本的比较，而不是模型额外输出了一个“犹豫”类别。
\end{example}

无误复核只是用于说明机制的简化条件。现实中人工也会出错，复核成本可能随队列长度变化，此时第三项必须包含这些后果。模型辅助决策的价值，应与可行替代方案在同一条件下比较，例如完全人工处理、简单规则或较低权限的自动化。只比较两个模型的准确率，会遗漏它们是否优于现有流程这一更基本的问题。

成本难以精确确定时，可以报告不同成本比下的决策变化。如果两种方案只在某个很窄的成本区间内有优势，使用者就能识别结论依赖了哪项判断；如果一个方案在相关区间中都更好，选择依据更稳固。这类敏感性分析检查的是使用假设的变化，应与第~\ref{chap:trust-robustness}章对输入和环境变化的分析区分。

预测损失与行动损失之间还可能隔着一个现实过程。若某个分数用于决定谁获得培训，接受培训本身会改变后续成绩；此时，成绩不再是完全独立于决策的固定标签。只在过去获得培训的人中测得预测准确，不能直接回答把名额分给其他人会产生什么收益。需要区分“预测现状”和“比较行动后果”，并说明后者所需的额外观测或因果假设。式~\eqref{eq:trust-system-risk}把$y$视为不随当前行动变化的状态，适合邮件真实类别这样的情形。若行动改变结果，就需要描述不同候选行动下的可能结果，再比较相应损失；不能直接把历史行动下观察到的成绩当成所有行动共用的$y$。第~\ref{chap:trust-fairness}章讨论干预的长期影响时，将再次遇到这一问题。

\subsection{风险的承担者与使用范围}

系统的购买者、操作人员和受影响者不一定是同一群体。招聘系统可能帮助用人方筛选简历，却主要由求职者承受误拒后果。评价时应同时记录谁获得收益、谁承担错误、谁能够发现和纠正问题。只有开发者方便测量的指标，未必覆盖这些不同影响。

风险还取决于使用范围。语言、设备、时间、输入质量和用户群体都可能改变模型的有效性。部署说明中的“适用于某类输入”，应能对应到可检查的条件；如果遇到范围外输入，系统需要有识别、限制使用或转交的安排。第~\ref{chap:trust-robustness}章将讨论环境变化如何使已有规律失效。

失效不只来自预测错误。系统可能预测正确，却暴露了私人信息；可能忠实完成用户指令，却执行了不被允许的动作；也可能被外部内容诱导而偏离任务。把这些后果写成具体情形，才能确定后续需要什么保护。NIST的人工智能风险管理框架同样强调结合使用情境识别、测量和管理风险，但框架中的原则仍需落实为当前系统的具体要求\scite{nist2023rmf}。

\section{系统属性与规范要求}
\label{sec:trust-requirements}

要求的设定包含相互依赖的三步：把后果写成可判断的主张，检查这些主张是否遗漏重要情境，再处理不同要求之间的支持或冲突。一个可检验的要求，需要把希望避免的后果、能够观察的事件与决定行动的规则连起来。只有后果描述而没有可观察事件，无法判断系统是否改进；只有一个测量指标而没有行动含义，也可能优化了与使用目的无关的数值。

\subsection{可信性要求的具体表达}

“模型很可靠”没有说明可靠的对象。更明确的主张可以是：对规定分布中的邮件，自动隔离部分的正常邮件比例不超过某个阈值；或者，对指定人群给出的预测区间具有至少某个覆盖率。每条主张都需要标明评价单位、总体范围、错误事件和容许水平。

评价单位尤其容易被忽略。以单次请求为单位的失败率，与以用户一周的全部请求为单位的失败率不同。一个每次调用都很少出错的工具，在被连续调用后仍可能产生较高的累计失败概率。多步任务应同时检查单步性质与最终结果，不能把前者直接当成后者。

要求也不都适合写成平均损失。有些行为需要限制发生概率，有些要求需要对不同群体分别成立，还有些属于过程约束，例如敏感数据的访问必须受到授权。可以把目标概括为在满足约束的前提下选择模型与决策规则：
\begin{equation}
 \min_{f,d}\ R_{\mathrm{sys}}(f,d),
 \qquad
 G_j(f,d)\leq b_j,\quad j=1,\ldots,m.
 \label{eq:trust-constrained-objective}
\end{equation}
这里$G_j$是第$j$项可量化要求，$b_j$是其容许水平，$m$为要求数量。无法合理量化的程序要求仍需单独保留，不能为了套用公式而虚构分数。

\subsection{可信性要求的完整性}

即使已经约定风险，训练与验证也未必唯一确定了我们希望得到的行为。\term{欠规定性}（underspecification）指已有任务定义和评价条件允许多个表现近似相同的模型，却没有约束它们在实际关切上的差异。D'Amour等在不同机器学习流程中展示了这一问题：训练域的留出表现接近，并不意味着部署时的行为可以互换。这将诊断对象从单个模型是否过拟合，扩展到整套训练和评价条件是否足以筛选模型\scite{damour2020underspecification}。

可以用一个构造例子看出二者的区别。设邮件训练数据中，内容特征$x_1$与页脚特征$x_2$总是相同，且标签等于二者。使用$x_1$或$x_2$的分类器在独立同分布的测试集上都能完全正确；继续增加同一来源的测试邮件，仍无法决定该信哪一条规律。如果部署后只有内容关系保持，两个模型便不再等价。这里缺少的是能区分候选规律的情境，单纯减少同一评价分布上的估计误差不能补上这类信息。

相应的改进也有两个层次。训练可以利用更多环境、结构约束或领域知识；评价则要主动设置能区分候选规律的测试。选择哪一层补充信息，取决于实际部署是否允许取得这些信息。不能因为某个随机种子碰巧在压力测试上表现较好，便把所有重新训练的模型都视为已通过同一验证。

另一种遗漏发生在类别内部。\term{隐藏分层}（hidden stratification）指数据中存在与使用后果有关的子类型，却没有在标签或评价分组中表示。Oakden-Rayner等的医学影像研究讨论了低频子类型、标注质量及伪相关怎样使总体表现掩盖重要子集的失败。其启示是检查现有标签是否描述了足够的任务差异，而不是认定所有自动聚类得到的组都具有实际意义\scite{oakdenrayner2019hidden}。

在邮件例子中，“正常邮件”可能同时包含普通通知和有截止日期的合同。即使两者的数量差异很大，使用者也可能更在意后者的漏读。应由任务后果提出候选切片，再核验每片的样本质量与统计不确定性。依据错误发现的新切片可以帮助诊断，但发现切片与确认切片需要区分，否则反复搜索后最差的一组可能只是抽样波动。鲁棒性与公平性两章将分别讨论环境切片和群体切片，本章先明确它们共同的出发点：总体指标不能代表未被表达的使用要求。

\subsection{不同性质之间的联系与取舍}

本部分讨论的性质对应不同问题。可靠性检查预测及其不确定性表达是否有依据；鲁棒性考察输入或环境变化后的有效性；公平性关注机会、错误与负担的分配；目标对齐与行为约束考察学习到的行为是否符合实际要求。可解释性提供理解和质疑这些行为的证据；隐私保护限制数据使用中的信息泄露；对抗安全进一步考察有人主动操纵系统时，这些边界能否保持。

这些性质可以相互支持，也可能产生取舍。例如，把更多输入转交人工可能降低自动处理部分的错误率，却减少自动处理比例；更严格的隐私保护可能增加统计估计的噪声；限制工具权限可能减少损害，也可能使任务无法完成。评价应报告同一使用目标下的变化，而不能只展示被优化的那一个指标。

将不同指标加权求和，还隐含了允许相互补偿的假设。如果未经授权披露数据属于不可接受行为，较高的准确率就不能自动抵消它。式~\eqref{eq:trust-constrained-objective}将约束与优化目标分开，正是为了明确这种区别。阈值和优先级需要由使用目的、受影响者与领域知识共同确定，学习算法只能在给定要求下求解。

\section{证据能支持什么结论}
\label{sec:trust-evidence}

要求明确后，还需区分证据能支持多强的结论、如何取得有效证据，以及如何保留复核条件。这三个问题分别对应证据的含义、产生和记录。系统改变后原证据是否仍然适用，则属于下一节讨论的生命周期条件。

为了避免只记录一个有利分数，可以用十项内容组织一项\term{可信保证}（trustworthiness assurance）。这里的“保证”不是无条件承诺，而是一份可以逐项核对的论证记录。它先说明情境与危害，写清任务、受影响对象、评价单位和行动链；随后提出要求，明确指标或程序约束的分母、阈值和适用范围。

证据部分需要继续交代显式假设、论证、多源证据，以及反证与剩余不确定性。显式假设覆盖抽样关系、标签、权限、资源和替代方案；论证说明这些条件为何能支持要求，并保留量词与比较基准；多源证据可以来自理论、随机抽样、压力测试、机制分析、人因评价和监测；反证与剩余不确定性则记录失败案例、未覆盖情境、估计误差与替代解释。

最后四项把技术结论接入实际决策与生命周期管理。责任主体说明谁准备证据、作出决定、监测并执行处置；风险接受记录谁在什么范围、期限和不可补偿约束下承担剩余风险；失效触发指出哪些数据、模型、权限、流程变化或监测结果会暂停沿用原结论；处置则规定降级、停用、回退、通知、补救、复验与重新开放。威胁条件不另列一项，相关信息分别进入情境与危害、显式假设和失效触发；存在主动攻击者时，还要写明其知识、权限、资源、目标以及可能形成的损害路径。

\subsection{理论保证与经验检验}

理论保证解释在特定假设下必然成立的性质。它的力量来自量词和条件，而不是“有证明”这一标签。例如，一个覆盖保证可能对校准数据和未来样本的联合抽样成立，却不保证每个固定输入都具有相同覆盖率。使用保证时，需要检查数据生成方式、算法实现和实际评价对象是否与定理一致。

经验检验观察具体样本与攻击条件下的表现。它能够发现理论假设与现实之间的差距，也能够比较没有完整理论刻画的方法，但未观察到失效不等于证明失效不可能发生。下面的计算说明，即使评价设计很简单，也需要区分这两种结论。

\begin{example}[零次失败的含义]
\label{ex:trust-zero-failures}
设固定系统面对$n$个独立同分布试验，每次失败概率为$p$。观察到零次失败的概率是$(1-p)^n$。令$(1-p)^n=\delta$，可得到零失败结果对应的单侧置信上限
\[
 p_{\mathrm{up}}=1-\delta^{1/n}.
\]
在95\%置信水平下，100次零失败对应的上限约为2.95\%，1000次零失败对应的上限约为0.30\%。这里的95\%描述构造上限的统计程序在重复抽样中的覆盖性质，不是给固定但未知的$p$赋予95\%的后验概率。这些计算要求试验单位与抽样假设成立；反复尝试直到出现一批零失败结果后再报告，不能继续沿用同一解释。
\end{example}

\subsection{检验设计与证据有效性}

\term{行为测试}（behavioral testing）将“模型应当如何响应某种输入变化”写成可以执行的检查。Ribeiro等提出的CheckList把语言能力与测试类型结合，使评价不只依赖随机测试集的总准确率。它所组织的最小功能、保持不变和定向变化等检查，帮助开发者把对行为的预期变成具体用例\scite{ribeiro2020checklist}。

沿用邮件任务，可以从这三类预期构造测试。最小功能检查给出简短且含义明确的例子，例如独立的会议通知；保持不变的检查只改动大小写或无关排版，要求正常邮件仍被放行；定向变化则在任务已将某类未经请求的推销认定为垃圾邮件的前提下，逐步加入明确的推销内容，检查垃圾邮件分数是否朝预期方向变化。每项测试都依赖任务语义。将“转发诈骗警示中的诈骗原文”一律当成诈骗，恰好会暴露模型没有区分引用与实际意图。

行为测试与随机抽样检验各有职责。前者可以有目的地暴露机制缺口，后者可以在抽样假设下估计某个总体的风险。一个压力测试集故意大量收集困难样本，其失败率一般不能直接当作部署失败率；反过来，在大规模普通样本上未发现某个低频缺陷，也不足以免除针对它的检查。只有说明样本怎样取得，数值才有相应解释。

模型选择、阈值调整和最终评价承担不同职责。如果根据测试结果反复修改模型，原来的测试集就逐渐成为开发数据。需要保留未参与选择的数据，或者使用能处理多次选择的统计方法。第~\ref{sec:reliability-guarantee}节将给出有限候选程序的一个简单保证。

比较两个系统时，应尽量保持任务、样本、资源预算和评价口径一致。对于随机生成系统，还需要记录生成次数、解码设置和选择规则。只挑选最好的回答与另一个系统的第一次回答比较，会把额外计算和选择机会混入所谓模型优势。

人工或模型评判者也属于测量过程。需要检查其判断标准、重复一致性及与实际任务后果的关系。一个偏爱长回答的评判者，可能提高冗长答案的得分，却没有提高事实正确率。评价证据应保存模型版本、数据来源、评价规则和必要的运行条件，使其他人能够复核结论。

\subsection{证据的记录与复核}

一个测试结果之所以能够被正确使用，往往依赖没有写在分数中的信息。测试数据是否包含重复记录、标签来自谁、哪些样本被排除、与预训练语料是否重合，都会影响结论。Gebru等提出的Datasheets for Datasets主张随数据提供动机、组成、收集过程与推荐用途等说明，使数据使用者能够判断它与当前任务是否相容。这是一种记录数据条件的方式，本身并不证明数据没有偏差\scite{gebru2021datasheets}。

模型也需要对应的说明。Mitchell等提出Model Cards，强调记录预期用途、评价过程以及相关群体或条件下的表现。它把一个笼统的模型名称具体化为具有版本、范围与证据的对象。模型卡可以承载本章所说的可信主张，但如果缺少真实评价，填写完整的卡片仍不能代替证据\scite{mitchell2019modelcards}。

二者的关系可以通过一次复核理解。邮件模型的记录称其正常邮件误拦率较低，但数据说明显示测试样本全部来自某一种语言和某个用户群体。复核者据此能发现，结论尚未覆盖另一种语言，而不是误把模型卡中的一个平均数理解为普遍承诺。如果之后补充了新的测试，更新的应当是具体覆盖范围和证据，不能只把模型版本号改掉。

记录还有助于区分不同复现目标。重新运行同一程序得到接近的数值，检查的是计算与随机性；用新的独立样本得到一致结论，检查的是跨独立样本的复现一致性；在新的环境中仍满足要求，检查的是外部适用性。三者扩大了复核范围，但相应结论并不相互推出。公开代码有助于第一项，对后两项仍需要合适数据及明确的环境联系。

对于生成系统，保存一个最终回答尤其不够。模型版本、系统提示、检索内容、工具权限、采样温度、重试次数与答案选择方式都可能改变输出分布。若比较时名义上只考察模型变化，实际却同时增加了重试机会，观测到的收益应归于整个过程，不能全部解释为模型能力提高。这些条件应随评价结果一起保存，也应出现在之后的重新验证中。

\section{威胁与生命周期治理条件}
\label{sec:trust-lifecycle}

前两层说明“希望系统做到什么”和“现有证据支持到哪里”，还没有回答这些结论怎样随系统一起使用。生命周期条件需要记录哪些参与者、环境变化或系统组件可能使结论失效，谁有权接受剩余风险，哪些观测触发限制使用，以及怎样恢复到可以再次评价的状态。本章说明审批、监测和复查需要哪些技术材料，但不替组织规定法定职责，也不替受影响者决定什么后果可以接受。

\subsection{威胁条件、责任主体与风险接受}

威胁条件不只指恶意攻击。数据来源变化、标注流程失真、供应商更新、权限扩大、人员负荷和组件故障，都可能破坏保证所依赖的条件。需要记录变化能够作用在哪个对象、通过什么路径影响要求，以及现有监测能否发现。有人主动选择这些变化以造成损害时，还要加入其知识、权限、资源和目标；这些条件共同构成分析主动攻击时的威胁模型。

\term{责任主体}（accountable role）是对保证中的某项决定或行动负有明确职责的角色。准备评价、独立复核、批准使用、运行监测和事故处置可以由不同角色承担；只写“项目组负责”会让触发后无人知道谁能停用系统。\term{风险接受}（risk acceptance）则是在已知证据边界和剩余不确定性的前提下，由有权主体决定某项残余风险是否能在规定范围与期限内承担。它不是证明风险消失，也不能让一个指标的收益抵消已经声明为不可接受的约束违反。

\subsection{变更触发与组合后的重新验证}

一个组件满足要求，并不意味着组件组合后仍满足同样要求。检索器返回的内容可能改变生成模型的输入分布，自动重试可能扩大隐私泄露机会，人工审核也可能因告警过多而失效。重新验证应围绕这些新关系展开，而不只是重复组件测试。

可信主张因而需要附带\term{失效触发}（invalidation trigger）：哪些数据变化、权限扩大、模型更新或监测结果会暂停沿用原结论。记录这些条件并不削弱结论，反而使使用者知道何时能够依赖它。组织如何审批、发布和追责属于更完整的治理流程；技术材料至少要暴露决策点、所需证据和触发条件，使这些流程能够作出有依据的判断。

系统组合还可能改变错误之间的依赖关系。假设在同一垃圾邮件总体上，两个检测器各有10\%的漏检率，并且任一个报警就能拦截邮件。只有在这一总体上两个漏检事件独立的条件下，才能把系统的共同漏检率算成1\%。若它们使用同一表征并遗漏同一类邮件，共同漏检率可能仍为10\%。设置第二个检测器的价值因此要看它是否提供新的信息，而不能只看它在结构上是否“多了一层”。

人工复核也需要同样处理。人看到的可能是原始证据，也可能只是模型生成的摘要；摘要已经遗漏关键信息时，增加确认步骤未必纠正错误。若要主张人机协作降低了风险，应评价真实的协作过程，包括复核者能看到什么、需要处理多少告警，以及是否有时间采取行动。第~\ref{sec:alignment-oversight}节将进一步把这些要求落实到行为监督。

\subsection{失效处置与恢复}

触发失效后，处置目标是先限制损害，再保存足以定位问题的证据，随后恢复到可以重新评价的状态。根据危害类型，可以暂停自动动作、降低权限、切换到人工流程、回退模型或隔离数据；已经删除的内容、已经泄露的信息和仍可撤销的状态具有不同恢复边界，不能统一写成“回滚即可”。处置还应说明如何通知受影响者、如何纠正已有结果，以及哪些日志必须在满足隐私与访问约束的条件下保留。

重新开放不是原版本再次启动，而是对触发原因和修复范围重新建立证据。复验应包含已知失败的回归检查，也要检查修复是否改变正常任务、其他群体或相邻性质。责任主体据此重新判断剩余风险；若证据仍不足，就应继续限制使用或选择危害更小的替代流程。监测、处置与恢复因此不是部署后的附加步骤，而是保证成立条件的一部分。

\section{例子：邮件隔离的完整保证记录}
\label{sec:trust-mail-assurance}

前三节已经分别说明如何设定要求、评价证据并管理结论的生命周期。下面沿用章首邮件系统，把这些要素合并为一个完整的教学验收实例。模型、样本、阈值和结果均为人工构造，只用于展示怎样形成可判定的保证记录，不能当作真实邮件服务的经验数据。实例固定模型版本$v_1$及其概率修正程序$c_1$，适用于公司自有网关在批准后30日内接收的中文或英文文本邮件；加密正文、纯图像邮件和从网关外导入的邮件不在范围内。误拦可能延误合同与求职通知，漏放可能使用户接触诈骗内容，系统因此不获得永久删除权限。规则在$p<0.2$时放行，在$0.2\leq p\leq0.98$时转人工复核，在$p>0.98$时进入可恢复隔离。

独立验收集包含600封正常邮件和400封垃圾邮件。规则隔离362封，其中2封为正常邮件；转人工100封，其中70封为正常邮件；放行538封，其中10封为垃圾邮件。验收要求隔离量不少于300封、自动隔离比例不低于30\%，且以隔离邮件为分母的正常邮件比例，其单侧95\%二项Clopper--Pearson置信上界不超过2\%。当前点估计为$2/362\approx0.55\%$，置信上界约为1.73\%，自动隔离比例为36.2\%，三项均达到规定阈值。50次恢复演练从用户点击“恢复”计时，直到邮件同时出现在收件箱和搜索索引中；第95百分位为3分钟，最大值为7分钟，分别低于5分钟和10分钟的限制。

这项结论显式假定验收集与30日适用范围具有约定的抽样关系，标签可靠，$v_1$、$c_1$和两个阈值在验收前固定，人工队列可用，隔离区和搜索索引均支持恢复。论证由四个环节组成：行动损失确定三段规则，范围检查限制自动处理对象，非零最低隔离量与置信上界控制抽样不确定性，恢复演练检验补救路径。除总体计数外，中英文各有500封邮件，每组隔离181封并误拦1封；40项预先规定的行为压力测试也全部通过。

这些证据仍留下明确缺口。10封垃圾邮件被放行，新式模板、隐藏子类型、共同训练盲点和人工复核错误仍可能使系统失效；低频重要邮件的估计区间可能较宽，目前也没有真实上线监测结果。评价人员负责冻结验收集并计算上界，产品责任人批准范围与有效期，运行人员负责监测、停用和回退，用户则能够查看隔离结果并发起恢复。根据上述结果，教学实例中的产品责任人接受$v_1$在规定范围内运行30日，只允许可恢复隔离，不接受永久删除；隐私或安全约束失败时，也不得由较高准确率抵消。

模型、修正程序、阈值、标签规则或工具权限一旦改变，邮件超出声明范围，隔离量低于300封、比例低于30\%、置信上界或恢复时延越界，人工队列不可用，或者出现攻击与泄露迹象，原结论就停止适用。此时应暂停自动隔离并转入人工流程，回退到已验证版本，恢复误拦邮件并保留事件记录；补充评价和处置证据，经责任主体重新接受后再逐步开放。

\section{本章小结}
\label{sec:trust-foundations-summary}

章首的邮件分类器不能仅凭准确率决定是否自动删除邮件。三个层次把这个问题拆开：系统属性或规范要求规定自动处理的范围、错误分母和不可补偿约束；证据能力说明独立样本、压力测试与机制检查能够支持什么；威胁与生命周期条件则规定由谁接受剩余风险，哪些变化触发停用，以及怎样恢复。前面的教学实例据此只支持有条件、可恢复的自动隔离，并不支持永久删除。其中2\%上限以被隔离邮件为分母，与以全部正常邮件为分母的误拦率不同；两者若都影响使用决定，就必须分别记录。

不同性质回答的问题并不相同。可靠性检查概率与风险承诺是否得到数据支持，并讨论何时把输入转交其他流程；鲁棒性处理规定变化下的有效性，公平性检查机会、错误与负担的分配；目标对齐与行为约束、可解释性、隐私保护和对抗安全分别补充行为边界、解释证据、信息保护与主动攻击条件。它们可以相互提供证据，却不能由某项改进替代其余要求。

技术可信性最终形成的是一组有边界、可复核且带处置条件的主张，而不是一个综合标签。技术证据可以揭示系统在哪些条件下达到要求、还留下哪些不确定性，并为审批和复查提供输入；它不能替代组织对职责的安排、法律判断或受影响者参与的公共治理决定。

\paragraph{延伸阅读}

NIST人工智能风险管理框架适合继续理解使用情境、测量和风险处置之间的关系；阅读时应把框架原则落实为具体系统中的责任主体、证据和失效触发，而不是只保留抽象分类\scite{nist2023rmf}。

若希望进一步建立可复核的技术记录，可对照数据集说明书与模型卡。前者强调数据的动机、组成和采集条件，后者强调模型的预期用途与分组评价；二者都帮助限定证据范围，却不能替代实际测试\scite{gebru2021datasheets,mitchell2019modelcards}。
